\documentclass[10pt,a4paper]{amsart}
\usepackage[margin=19mm]{geometry}
\usepackage{amsmath,amssymb,mathtools}
\usepackage[scr=rsfs,cal=euler]{mathalfa}
\usepackage{xfrac}
\usepackage[unicode,hidelinks,bookmarks=false]{hyperref}
\newcommand{\ie}{i.e.\ }
\newcommand{\cf}{cf.\ }
\newcommand{\resp}{respectively\ }
\newcommand{\Subsection}{\subsection}
\providecommand{\nobreakdash}{\nobreak\hbox{-}}
\setlength{\emergencystretch}{2em}
\multlinegap0pt
\usepackage{amssymb}
\usepackage[shortlabels]{enumitem}
\newtheorem{remark}{Remark}
\newtheorem{lemma}{Lemma}
\newtheorem{theorem}{Theorem}
\title{Iwahori Spherical Whittaker Functions for Steinberg Representations}
\author{Markos Karameris, Ehud Moshe Baruch}
\date{Source: arXiv:2407.01448v6 (CC BY 4.0)}
\begin{document}
\maketitle

\noindent\emph{Source and licence.} Iwahori Spherical Whittaker Functions for Steinberg Representations. Version arXiv:2407.01448v6 (CC BY 4.0), distributed by arXiv under the Creative Commons Attribution 4.0 International licence, http://creativecommons.org/licenses/by/4.0/, as recorded in the arXiv licence metadata for this exact version and shown on the abstract page https://arxiv.org/abs/2407.01448v6. Full licence notice: \texttt{licenses/arxiv-2407.01448-CC-BY-4.0.txt}.\par\smallskip

\noindent\textit{Editorial scope.} Counted unit: the article's theorem main (source0.tex lines 191--197). The article's complete original proof of it is delivered with it (source0.tex lines 198--214, one \texttt{proof} environment of 17 lines). No mathematics is written, repaired or re-derived: the statement and the proof are the article's own lines, in the article's own notation, and no retained line is substituted or re-flowed.  Retained uncounted as the source-local context the proof needs: lines 129--134; lines 145--148; lines 161--168; lines 176--178; lines 183--185.  Nothing was omitted from any retained span, so the package carries no omission record.  No retained line contains a citation, so the wrapper carries no bibliography and no bibliography entry.  No retained line contains a figure environment, an image-inclusion command or drawing code, so no image is delivered and the package declares no asset dependency.  Editorial formatting only, in addition: an AMS \texttt{amsart} preamble that restates the article's own declarations and macros used by the retained lines, namely the environments \texttt{remark}, \texttt{lemma}, \texttt{theorem} on separate counters, together with the standard packages those lines need.  The retained environments print in this document as Remark 1, Lemma 1, Theorem 1 in order of appearance, whereas the article prints the same items with the numbering of its own full document.  The complete native source of the article as distributed in the arXiv e-print for this exact version is bundled unchanged as \texttt{sources/161641/source0.tex}.
\begin{remark} \label{domd} It is straightforward to verify that $d_w$ is the unique element in $T/T^0$ satisfying the condition
 $\langle\alpha,\lambda_{d_w}\rangle= \begin{cases}
        0 & \text{if $w^{-1}\alpha\in \Phi^+$} \\
        -1 & \text{if $w^{-1}\alpha\in \Phi^-$} 
    \end{cases},\forall\alpha\in\Delta$.
\end{remark}

\begin{lemma} \label{permlem}
         For a Weyl element $w$ and $d\in T^+$ the following relation is satisfied:
 $$W(dw)=(-q)^{-\ell(w)}W(d).$$
\end{lemma}

\begin{theorem}[Diagonal Whittaker values] \label{diagon} ~\\
 Let $W\in \mathcal{W}(\pi_{St^{\chi}},\psi)^{J}$ with $W(1)=1$ be such that $W$ is an eigenfunction of $\mathcal{H}(G,J)$ with $$(F*W)(g)=\rho_{\chi}(F)W(g).$$
	Then $W$ has the following diagonal values:
$$W(d)=\begin{cases}
       (\chi\delta_B)(d)  & \text{if $d\in T^+$} \\
        0 & \text{otherwise} 
    \end{cases}.$$
 \end{theorem}

\begin{lemma} \label{diagcos}
  For any $w\in\mathbf{W}$ we have $Jd_wwJ=\underset{\ell\in L_w}{\bigcup}n_{\ell}^-d_wwJ$, where $n_{\ell}^-\in N^-_{\mathfrak{p}}$, $L_w$ is a finite set of indices and $|L_w|=\mu(Jd_wwJ)$.
\end{lemma}

\begin{lemma}\label{aux}
    For every Weyl element $w\in\mathbf{W}$ we have $d_{w_0w}=w_0d_ww_0d_{w_0}z_w$ for some $z_w\in\mathcal{Z}$.
\end{lemma}

\begin{theorem}[Whittaker values at every cell $Tw$] \label{main} ~\\
Under the assumptions of Theorem \ref{diagon}
    $$W(dw)=\begin{cases}
      (\chi\delta_B)(d)(-q)^{-\ell(w)}  & \text{if $\lambda_d$ is $w$-dominant} \\
        0 & \text{otherwise} 
    \end{cases}$$
\end{theorem}

\begin{proof}
We use again the relation $X_{d_ww}(W)(g)=\epsilon_{d_ww}W(g)$. By Lemma \ref{diagcos} we have
$\epsilon_{d_ww}W(g)=\sum\limits_{\ell\in L_w}W(gn_{\ell}^-d_ww)$. We set $g=dd_{w_0}w_0(=dw_0d_{w_0}^{-1})$ and write $n_{\ell}^-=\prod\limits_{\alpha\in \Phi^-}x_{\alpha}(\varpi s_{\alpha})$ for some $s_{\alpha}\in\mathcal{O}$ depending on the indice $\ell$. We want to show that $\psi(gn_{\ell}^-g^{-1})=\psi(dw_0d_{w_0}^{-1}n_{\ell}^-d_{w_0}w_0d^{-1})=1$ for any $d\in T^+$. For this purpose we distinguish two cases for each $x_{\alpha}(\varpi s_{\alpha})$: 
\begin{itemize}
    \item for any $\alpha\in\Delta^-$: $w_0d_{w_0}^{-1}x_{\alpha}(\varpi s_{\alpha})d_{w_0}w_0=x_{w_0(\alpha)}(\varpi s_{\alpha}\varpi^{\langle\alpha,\lambda_{d_{w_0}^{-1}}\rangle})\in N_{\mathcal{O}}$ since $\langle\alpha,-\lambda_{d_{w_0}}\rangle={-1}$ by Remark \ref{domd}. Setting $n^+=w_0d_{w_0}^{-1}x_{\alpha}(\varpi s_{\alpha})d_{w_0}w_0$, it follows that $\psi(dn^+d^{-1})=1$ for any $d\in T^+$.
    \item for any $\alpha\in\Phi^-\backslash\Delta^-$ and $d\in T$: $\psi(dx_{w_0(\alpha)}(\varpi s_{\alpha}\varpi^{\langle\alpha,\lambda_{d_{w_0}^{-1}}\rangle})d^{-1})=1$ since $w_0\alpha\not\in\Delta$.
\end{itemize} For $g=dd_{w_0}w_0$ we then have that $W(gn_{\ell}^-g^{-1}gd_ww)=W(gd_ww)$, giving the relation:  $$\epsilon_{d_ww}W(dd_{w_0}w_0)=\mu(Jd_wwJ)W(dd_{w_0}w_0d_ww)=\mu(Jd_wwJ)W(dd_{w_0}w_0d_ww_0w_0w)$$ for any $d\in T^+$. Using Lemma \ref{aux} now, this expression becomes:
$$\epsilon_{d_ww}W(dd_{w_0}w_0)=\mu(Jd_wwJ)W(dd_{w_0w}w_0w),\text{ for all } d\in T.$$
Notice that the above relation is extended to all of $T$ and not just $d\in T^+$ since $d\not\in T^+$ means $dd_{w_0}\not\in T_{w_0}^+$ and $dd_{w_0w}\not\in T_{w_0w}^+$, so both sides are $0$.
Setting $w=w_0$ and comparing scalar terms with Lemma \ref{permlem} for $d=d_{w_0}^{-1}$ we get that $\epsilon_{d_{w_0}w_0}=(-q)^{\ell(w_0)}\mu(Jd_{w_0}w_0J)W(d_{w_0}^{-1})$. Substituting the value of $\epsilon_{d_{w_0}w_0}$ and setting $d\to dd_{w_0}^{-1}$ leads to: 
$$W(dw_0)=(-q)^{-\ell(w_0)}W(d_{w_0}^{-1})^{-1}W(dd_{w_0}^{-1}), \text{ for all } d\in T$$ Similarly we set $w\to w_0w$ and compute $\epsilon_{d_{w_0w}w_0w}=(-q)^{\ell(w_0)-\ell(w)}\mu(Jd_{ww_0}ww_0J)W(d_{w_0}^{-1}d_w)$. Substituting this value of $\epsilon_{d_{w_0w}w_0w}$ and using the above relation for the $W(dw_0)$ term implies:
$$W(dw)=(-q)^{-\ell(w)}W(d_{w}^{-1})^{-1}W(dd_{w}^{-1}), \text{ for all } d\in T \; \; (**)$$ Note that $d_w^{-1}\in T^+$ so a direct application of Theorem \ref{diagon} proves the claim since: 
\begin{itemize}
    \item if $d\not\in T_w^+$ both sides evaluate to zero,
    \item if $d\in T_w^+$ then it follows that $dd_w^{-1}\in T^+$, and the expression evaluates to $W(dw)=(-q)^{-\ell(w)}(\chi\delta_B)(d_w^{-1})^{-1}(\chi\delta_B)(dd_w^{-1})=(-q)^{-\ell(w)}(\chi\delta_B)(d)$.
\end{itemize}
\end{proof}

\end{document}
